\begin{abstract}
We characterize the sharp expected length of honest confidence intervals for a homogeneous conditional log-odds coefficient with binary treatment and outcome. The experiment has a known uniform scalar covariate design, an effect bounded in absolute value by \(1/2\), and propensity and control-risk logits with supremum norms at most \(1\) and Hölder seminorms at most \(2\). For public propensity and prognosis smoothness exponents \(\alpha,\beta\) satisfying \(0<\beta<1/4\) and \(\beta<\alpha\le1\), coverage is required uniformly over this ambient model, while expected length is evaluated on laws with effect magnitude at most \(r\), the evaluation radius. For every sample size \(n\ge1\) and \(r\in[0,1/2]\), the minimax expected length is, up to absolute constants,
\[
n^{-\min\{1/2,\,2(\alpha+\beta)/(2\alpha+2\beta+1)\}}
+r\,n^{-4\beta/(4\beta+1)}.
\]
One radius-independent interval attains this profile simultaneously across all radii with finite-sample \(90\%\) coverage. Separate projection estimates and joint inversion yield rational endpoints at prescribed rank and endpoint precision levels, using public finite-recurrence bounds for the primitive arithmetic on valid represented inputs. Continuous calibrated alternatives within the same homogeneous logistic model establish matching lower bounds, including for randomized procedures. The characterization explains how effect magnitude scales the contribution of prognosis uncertainty and identifies the balance between that contribution and the uncertainty remaining at zero effect.
\end{abstract}

\section{Introduction}

This paper establishes the sharp expected-length profile for honest inference on a homogeneous conditional log-odds coefficient in a prescribed binary treatment and outcome experiment. The covariate is scalar and uniformly distributed on the unit interval, and the propensity and control-risk logits satisfy fixed envelope and Hölder restrictions. Their smoothness exponents are public inputs. Within this model, we construct one interval sequence with finite-sample \(90\%\) coverage and characterize its optimal expected length simultaneously across effect-radius subclasses. The characterization includes evaluation at zero effect, arbitrary shrinking-radius sequences, and fixed positive effect radii.

The target is \(\theta(P)\), the common conditional log-odds coefficient under an observed law \(P\). It is the treated-arm log odds minus the control-arm log odds at every covariate value. Under consistency and conditional exchangeability, this coefficient is also the homogeneous conditional causal log-odds contrast; \cref{obj:lem:observable-odds} records the corresponding full-data interpretation. Conditional odds contrasts have a longstanding role in stratified association analysis \citep{Mantel1959}, and potential-outcome identification supplies their causal interpretation \citep{Rubin1974,Rosenbaum1983}. Their conditional meaning is particularly relevant because confounding and odds-ratio collapsibility are distinct issues \citep{Greenland1999}.

Our precision criterion separates the coverage obligation from the laws at which interval length is evaluated. The ambient model in \cref{obj:def:model} imposes a known uniform design, a homogeneous effect with absolute value at most \(1/2\), nuisance-logit supremum norms at most \(1\), and Hölder seminorms at most \(2\). Let \(\alpha\) denote the public propensity-logit exponent and \(\beta\) the public control-risk-logit exponent. We study
\[
0<\beta<1/4,\qquad \beta<\alpha\le1.
\]
Every admissible procedure must cover the coefficient with probability at least \(90\%\) throughout this ambient model. For \(r\in[0,1/2]\), the effect evaluation radius, its worst-case expected length is assessed over the subclass with \(|\theta(P)|\le r\). This ambient-coverage and subclass-length distinction follows the honest-confidence-interval perspective of \citet{Low1997} and \citet[Section 2, Theorem 1]{CaiLow2004ConfidenceAdaptation}.

Write \(J_n(\alpha,\beta;r)\) for the minimax expected interval length under this criterion, with \(n\) the sample size. Our main result, \cref{obj:thm:full-honest-frontier-answer}, establishes
\[
J_n(\alpha,\beta;r)
\asymp
n^{-\mathsf a(\alpha,\beta)}
+r\,n^{-\mathsf b(\beta)},
\]
where \(\mathsf a\) and \(\mathsf b\), the numerator and denominator rate exponents, are
\[
\mathsf a(\alpha,\beta)
=\min\left\{\frac12,\frac{2(\alpha+\beta)}{2\alpha+2\beta+1}\right\},
\qquad
\mathsf b(\beta)=\frac{4\beta}{4\beta+1}.
\]
The positive comparison constants are absolute and apply simultaneously to every admissible exponent pair, every integer \(n\ge1\), and every evaluation radius. A single procedure, tuned by the sample size and public exponents, attains the upper comparison across all radii. The lower comparison applies to the entire ambiently honest class, including procedures with independent randomization.

The two terms have an observable-coordinate interpretation. The odds identity in \cref{obj:lem:observable-odds} expresses the coefficient through a numerator and a uniformly positive denominator. Numerator uncertainty contributes directly to interval length, whereas denominator uncertainty is scaled by effect magnitude. Separate projection resolutions and joint inversion turn this identity into the attaining interval of \cref{obj:def:upper-handle}. At zero radius, the sharp length is \(n^{-\mathsf a(\alpha,\beta)}\), attaining the parametric order precisely when \(\alpha+\beta\ge1/2\). Throughout the stated exponent domain, \(\mathsf a>\mathsf b\), so the denominator contribution eventually dominates at each fixed positive radius. For shrinking radii, the two contributions balance at
\[
r\asymp n^{-\{\mathsf a(\alpha,\beta)-\mathsf b(\beta)\}}.
\]
These comparisons describe effect-dependent precision while retaining the same ambient coverage requirement.

The construction also specifies a finite rational realization. An explicit interval-arithmetic engine and a dyadic naming policy select a concrete interval with ordered rational endpoints in the effect envelope on every sample. Under the quantitative primitive and representation contracts in \cref{obj:def:arithmetic-engine,obj:def:dyadic-name-convention}, every admissible engine and Borel naming policy satisfies the same coverage and expected-length comparisons. Public exponent representations are fixed before sampling. Prescribed rank and endpoint precision budgets establish termination on valid represented inputs; the resulting intervals may vary across admissible implementations.

The result connects estimand-specific logistic inference to projection-based functional estimation and honest interval optimality. \citet[Propositions 1--3]{Tan2019LogisticDR} develops robust logistic coefficient scores and effect-dependent efficiency comparisons near zero. \citet[arXiv version 1, Assumptions REG and ML1, Theorem 2]{LiuZhangZhou2020LogisticDML} establishes regular asymptotic inference under its nuisance-estimation and regularity conditions. Binary DINA likewise targets conditional log-odds contrasts \citep[Section 2.1.1, equation (5)]{GaoHastie2025DINA}. Projection-based conditional covariance estimation provides a closely related numerator construction \citep[version 5, Theorem 1 and Corollary 1]{McGrath2026}. Here, the contribution is the simultaneous sharp expected-length characterization under ambient honesty in the specified homogeneous logistic model. Its converse uses calibrated continuous alternatives whose every draw satisfies that model's homogeneity, envelope, and native-logit smoothness restrictions.

The results supporting \cref{obj:thm:full-honest-frontier-answer} hold under their stated assumptions.

The paper first reviews related work, then specifies the model and honest interval criterion in \cref{obj:def:model,obj:def:length-objective}. It next presents the sharp comparison in \cref{obj:thm:full-honest-frontier-answer}, develops the interval construction in \cref{obj:def:upper-handle}, and discusses interpretation and limitations. The appendices supply identification, projection bounds, rational realization, calibrated alternatives, original-record lower comparisons, the verification note, and proofs.
